\chapter{What the software produces}

\eyebrow{the deliverables}

\vspace{4pt}
The outputs of Apex are directed at a single purpose: a research dataset with
verifiable pairing, enriched with objective measurements. This chapter summarises
the outputs and the degree of confidence attached to each.

\section{The dataset}

The primary product is the \textbf{paired dataset}: for every image, a verifiable
link to its carcass, animal, treatment and stratum, plus the physical reference
measured on the animal. It is stored in the single \tele{carcass.db} and exported
as \tele{manifest.csv} + \tele{integrity\_report.txt}.

\section{Per-image measurements}

\begin{center}
\begin{tabularx}{\linewidth}{@{}l l X@{}}
\toprule
\textbf{Measurement} & \textbf{Status} & \textbf{What it is} \\
\midrule
Fat fraction \& mask & \pillok{validated} & Fraction of the carcass surface
      segmented as fat, with a cyan overlay. IoU $\sim$0.92. \\
Conformation convexity & \pillok{objective} & Per-region convexity (leg / loin /
      shoulder) and a convexity map. Analytical, invariant, training-free. \\
Finishing grade & \pillalert{experimental} & A 3-class finishing estimate from a
      model trained on n=22. Not validated. \\
Conformation grade & \pillalert{estimate} & A 5-class estimate derived from leg
      convexity. Not validated. \\
Fat-thickness (EG) & \pillalert{experimental} & A regression estimate. Not
      validated. \\
\bottomrule
\end{tabularx}
\end{center}

\section{Per-batch products}

\begin{itemize}[leftmargin=1.3em,itemsep=3pt]
  \item \textbf{Fat statistics}: mean, std dev, median, min, max, and the
        distribution histogram.
  \item \textbf{Inter-rater agreement}: Fleiss' $\kappa$ per axis and a
        consensus grade.
  \item \textbf{Sample progress}: coverage against the 100--120 target, by
        stratum.
  \item \textbf{CSV export}: one row per image with every measurement.
\end{itemize}

\section{Live products}

\begin{itemize}[leftmargin=1.3em,itemsep=3pt]
  \item \textbf{Live fat overlay}: real-time framing guide from camera or
        video (approximate).
\end{itemize}

\begin{note}[title=Reading the colour code]
Throughout the software, colour indicates the confidence attached to a number.
\textbf{Green} = validated or objective. \textbf{Amber} = experimental or an
estimate, to be validated later. \textbf{Red} is reserved for broken pairing and
is never used for a data series. Every coloured value carries a label, so colour is
never the only signal.
\end{note}
